%%%%%%%%%%%%%%%%%%%%%%%%%%%%%%%%%%%%%%%%%%%%%%%%%%%%%%%%%%%%%%%%%%%%%%
% Lớp 10 | Xác suất | Không gian mẫu và biến cố
% Dạng 0D0?1-1: Các câu hỏi lý thuyết tổng hợp
% Số câu: 5  (sinh tự động bởi tach_cau.py; từ nay thầy sửa trực tiếp file này)
%%%%%%%%%%%%%%%%%%%%%%%%%%%%%%%%%%%%%%%%%%%%%%%%%%%%%%%%%%%%%%%%%%%%%%

%%%%%%%%%%%%%%%% TRẮC NGHIỆM 4 PHƯƠNG ÁN %%%%%%%%%%%%%%%%

% nguồn: MaHoaNH_TracNghiem_Lop10
\begin{ex}%[0D0N1-1]%[ma:e4ad1a6f]
	Biến cố chắc  chắn kí hiệu là
	\choice 
		 {$A$}
		{\True  $\Omega$}
		{$\varnothing$} 
		 {$n\left(\Omega\right)$ }
	\loigiai{ 
		Biến cố chắc  chắn là biến cố luôn xảy ra, kí hiệu là $\Omega$.
	 }
\end{ex}

% nguồn: MaHoaNH_TracNghiem_Lop10
\begin{ex}%[0D0N1-1]%[ma:5199922e]
	Gieo ngẫu nhiên cùng lúc  $2$ đồng xu phân biệt cân đối và đồng chất. Số phần tử của không gian mẫu là
	\choice 
		{$2$} 
		{\True $4$} 
		 {$25$}
		 {$36$}
	 \loigiai{
		Trong một lần gieo, mỗi đồng  xu có $2$ khả năng xuất hiện.\\
	 Vậy số phần tử của không gian mẫu trong phép thử trên là $n(\Omega)=2\cdot 2=4$.}
\end{ex}

% nguồn: MaHoaNH_TracNghiem_Lop10
\begin{ex}%[0D0N1-1]%[ma:0762f298]
	Trong các thí nghiệm sau thí nghiệm nào  không phải là phép thử ngẫu nhiên.
	\choice 
		{Gieo đồng tiền xem xuất hiện mặt ngửa hay  mặt sấp}
		{Gieo 3 đồng tiền  và xem có mấy đồng tiền lật ngửa}
		{Gieo một con  xúc xắc và xem nó xuất hiện mặt mấy chấm}
		{\True Bỏ hai viên bi xanh và ba viên bi đỏ trong một chiếc hộp, sau đó lấy  từng viên một để đếm xem có tất cả bao nhiêu viên bi}
	 \loigiai{
		Phép thử ngẫu nhiên (gọi tắt là phép thử) là một hoạt động mà ta không thể biết  trước được kết quả của nó.\\
		Phương án {\bf D} không phải là phép thử vì ta biết chắc chắn kết quả chỉ có thể là một  số cụ thể số bi xanh và số bi đỏ.
	 }
\end{ex}


%%%%%%%%%%%%%%%% ĐÚNG/SAI %%%%%%%%%%%%%%%%

% nguồn: MaHoaNH_DungSai_Lop10
\begin{ex}%[0D0N1-1]%[ma:c2f33cab]
	 Chọn các khẳng định đúng trong các khẳng định sau
	\choiceTF 
		{\True ``Gieo $1$ con súc sắc'' có không gian mẫu là $\Omega=\{1; 2;  3; 4; 5; 6\}$}
		{Phép thử ``Gieo $1$  con súc sắc, sau đó giao tiếp $1$ đồng xu'', không gian mẫu là $\Omega=\{1; 2; 3; 4; 5; 6; S; N; S; N\}$}
		{\True ``Gieo  hai đồng xu phân biệt''. Nếu kí hiệu $S$ để chỉ đồng xu "sấp", kí hiệu $N$ để chỉ đồng xu "ngửa" thì không gian mẫu của phép thử trên là $\Omega=\{SS;SN;NS;NN\}$}
		{``Gieo $2$ con súc  sắc'' có số phần tử không gian mẫu là $n\left(\Omega\right)=12$}
	\loigiai{ 
		\begin{itemchoice}
			\itemch ``Gieo $1$ con súc sắc'' có không gian mẫu  là $\Omega=\{1; 2; 3; 4; 5; 6\}$ là khẳng định đúng.
			\itemch Phép thử ``Gieo $1$ con  súc sắc, sau đó giao tiếp $1$ đồng xu'', không gian mẫu là $\Omega=\{1; 2; 3; 4; 5; 6; S; N; S; N\}$ là khẳng định sai vì không gian mẫu $\Omega=\{1S;2S;3S,4S;5S;6S;1N;2N;3N,4N;5N;6N\}$.
			\itemch ``Gieo hai đồng xu phân biệt'' thì không gian mẫu của phép thử trên  là $\Omega=\{SS;SN;NS;NN\}$ là khẳng định đúng.
			\itemch ``Gieo $2$ con súc sắc'' có số phần tử  không gian mẫu là $n\left(\Omega\right)=12$ là sai vì $n\left(\Omega\right)=6^2=36$.
		\end{itemchoice}
	 }
\end{ex}

% nguồn: MaHoaNH_DungSai_Lop10
\begin{ex}%[0D0N1-1]%[ma:71960e1a]
	 Chọn các khẳng định đúng trong các khẳng định sau
	\choiceTF 
		{\True Phép thử ngẫu nhiên (gọi tắt là phép thử) là một thí nghiệm hay một hành động mà kết quả của nó không đoán trước được nhưng có thể xác định  được tập hợp tất cả các kết quả có thể xảy ra của phép thử đó}
		 {Phép thử ngẫu nhiên (gọi tắt là phép thử) là một thí nghiệm hay một hành động mà kết quả của nó có thể đoán trước được và có thể xác định được tập hợp tất cả các kết quả có thể xảy ra của phép thử đó}
		{\True Tập hợp mọi kết quả của một phép thử $T$ được gọi là \textbf{không gian mẫu}} 
		{Tập hợp mọi kết quả của một phép thử $T$ được gọi  là số phần tử của \textbf{không gian mẫu}}
	 \loigiai{
		\begin{itemchoice}
			\itemch Phép thử ngẫu nhiên (gọi tắt là phép thử) là một thí nghiệm hay một hành động mà kết quả của nó không đoán trước được nhưng có thể xác định được tập hợp tất cả các  kết quả có thể xảy ra của phép thử đó là khẳng định đúng.
			\itemch Phép thử ngẫu nhiên (gọi tắt là phép thử) là một  thí nghiệm hay một hành động mà kết quả của nó có thể đoán trước được và có thể xác định được tập hợp tất cả các kết quả có thể xảy ra của phép thử đó là sai vì phép thử không thể đoán trước được kết quả.
			\itemch Tập hợp mọi kết quả của một phép thử  $T$ được gọi là \textbf{không gian mẫu} là khẳng định đúng.
			\itemch Tập hợp mọi kết quả của một phép thử $T$ được gọi là  số phần tử của \textbf{không gian mẫu} là sai vì tập hợp mọi kết quả của một phép thử $T$ được gọi là \textbf{không gian mẫu}.
		\end{itemchoice}
	} 
\end{ex}
